% !TeX root = ../main.tex

\section{Objetivos}
\label{sec:objetivos}

\subsection{Objetivo general}
\label{subsec:objetivo-general}

Analizar la influencia de la variabilidad de los tiempos de servicio en el tiempo medio de espera en cola de un sistema M/G/1 con disciplina FCFS, manteniendo constantes la tasa de llegada y el tiempo medio de atención, mediante una revisión bibliográfica, una comparación analítica y una demostración con simulación de eventos discretos.

\subsection{Objetivos específicos}
\label{subsec:objetivos-especificos}

\begin{enumerate}
  \item Describir los fundamentos de los sistemas de colas con un servidor, identificando los supuestos de los modelos estudiados y distinguiendo la media y la variabilidad del tiempo de servicio, la utilización del servidor y el tiempo de espera en cola.

  \item Explicar, a partir de la fórmula de Pollaczek--Khinchine, la relación entre la variabilidad del tiempo de servicio y la espera media, y comparar los casos M/D/1 y M/M/1 como particularizaciones del modelo M/G/1 bajo iguales condiciones de llegada y media de servicio.

  \item Comparar críticamente los aportes de la bibliografía seleccionada sobre variabilidad y desempeño de sistemas de colas, diferenciando los resultados según sus supuestos, disciplinas de atención y medidas de desempeño.

  \item Desarrollar una demostración reproducible mediante simulación de eventos discretos que compare distribuciones de servicio con igual media y distinta variabilidad, y contrastar las estimaciones del tiempo medio de espera con los valores analíticos de referencia, considerando la incertidumbre estadística de las estimaciones.

  \item Interpretar los resultados de la comparación e identificar sus implicaciones y limitaciones, delimitando las condiciones bajo las cuales las conclusiones pueden aplicarse a los sistemas estudiados.
\end{enumerate}
